% !TeX root = ../教材题库_全部条目_转写工作稿.tex
% Source: XJ-B1-C04-HW-P147-Y1062
% 第4章 / 4.5 函数模型及其应用 / 教材 p.147、148 / PDF 第154、155页
% OCR未校对：公式与文字均可能有误；未完成数学TeX转写；禁止直接发给学生
% 原题及相邻原解截图（核对用）：../题目原图/XJ-B1-C04-HW-P147-Y1062-01.webp, ../题目原图/XJ-B1-C04-HW-P147-Y1062-02.webp
% 本文件只在转写完成后作为题干使用；OCR 条目尚不是可靠题干。
\begin{bankitem}{习题4.5 第3题}{第4章 / 4.5 函数模型及其应用 / 教材 p.147、148 / PDF 第154、155页}{XJ-B1-C04-HW-P147-Y1062}{OCR未校对：公式与文字均可能有误；未完成数学TeX转写；禁止直接发给学生}
\begin{ocrdraft}
[教材p.147 / PDF第154页]
\par
3．幂函数的增长快慢和幂指数的大小密切相关．但是，增长很快的幂函数和
\par
增长比较慢的指数函数相比，仍然是小巫见大巫．
\par
试用计算器计算并填写下表，探索这个现象．
\par
0
\par
1
\par
30
\par
50
\par
100
\par
巧0
\par
200
\par
250
\par
[教材p.148 / PDF第155页]
\end{ocrdraft}
\end{bankitem}
